\documentclass[11pt]{article}
\usepackage[margin=1in]{geometry}
\usepackage{amsmath,amssymb}
\usepackage{graphicx}
\usepackage[hidelinks]{hyperref}

\title{Literature Status: Real Two-Dimensional \texorpdfstring{$\ell_p$}{ell-p} Numerical Index}
\author{FA/Banach run \texttt{fa\_banach\_001}}
\date{2026-06-13}

\begin{document}
\maketitle

\section*{Status}

This is a literature-already-answered packet, not an original proof from the
run.  Problem 2 in Kadets--Martin--Paya, arXiv:math/0605781, asks whether
the numerical index of the real two-dimensional space $\ell_p^2$ is
\[
  \sup_{t\in[0,1]} \frac{|t^{p-1}-t|}{1+t^p}
  \qquad (1<p<\infty).
\]
Chiclana, arXiv:2606.01675, answers this affirmatively by proving that
$n(\ell_p^2)=v(J)$ for every $p\geq 1$, where
\[
  J=\begin{pmatrix}0&1\\ -1&0\end{pmatrix}.
\]

\section*{Original Problem}

The original source is:

\begin{quote}
V. Kadets, M. Martin, R. Paya, \emph{Recent progress and open questions on
the numerical index of Banach spaces}, arXiv:math/0605781.
\end{quote}

The paper states that, in the real case, the displayed supremum above is the
numerical radius of the norm-one operator $J$ on $\ell_p^2$, and asks whether
this upper bound is always the numerical index.

\begin{center}
\includegraphics[width=\textwidth]{figures/open_problem_crop.png}
\end{center}

\section*{Later Answer}

The supporting source is:

\begin{quote}
R. Chiclana, \emph{The Numerical Index of Two-Dimensional Real
$\ell_p$ Spaces}, arXiv:2606.01675.
\end{quote}

The abstract states that the paper computes the numerical index of
two-dimensional real $\ell_p$-spaces for all $p\geq 1$ and confirms the
conjectured formula in the two-dimensional real case.  Theorem 1.1 states:
\[
  n(\ell_p^2)=v(J)\qquad (p\geq 1).
\]

\begin{center}
\includegraphics[width=\textwidth]{figures/supporting_answer_crop.png}
\end{center}

Combining this theorem with the original source's identification of
$v(J)$ with the displayed supremum gives the affirmative answer to Problem 2
of arXiv:math/0605781 for all $1<p<\infty$.

\section*{Verification Notes}

The supporting paper is a separate later paper.  It identifies the same
two-dimensional real $\ell_p$ conjecture by formula and proves it.  In the
text inspected here, the supporting paper does not appear to cite
arXiv:math/0605781 directly; it cites later partial-progress papers and
describes the same conjecture.

The packet does not address the neighboring 2006 questions about positivity
of $n(\ell_p)$, equality $n(\ell_p)=n(\ell_p^2)$, or the complex
two-dimensional formula.

\section*{Search Bounds}

The registry and solution index were checked for arXiv:math/0605781,
arXiv:2606.01675, Chiclana, and the real two-dimensional $\ell_p$ formula.
The original and supporting arXiv PDFs were downloaded and inspected.  No
broader web, MathSciNet, or Zentralblatt search was performed.

\section*{Review Recommendation}

Verify the two displayed links in the transfer: the original source's
identity between the supremum and $v(J)$, and the supporting theorem
$n(\ell_p^2)=v(J)$.  No other mathematical input is used.

\end{document}
